\section{Introduction}
\label{sec:intro}

Looped Transformers, first introduced as Universal Transformers \citep{dehghani2018universal}, have recently re-emerged as a compelling alternative to conventional depth scaling. Rather than stacking distinct layers, they repeatedly apply the same model across recurrent steps, closely connecting this approach to parameter sharing in Transformers \citep{lan2019albert,dehghani2018universal}. This recurrent approach has shown strong empirical performance across a wide range of domains, including language modeling, algorithmic learning, and abstract reasoning \citep{gao2025universal,yang2023looped,jolicoeurmartineau2025trm,wang2025hierarchical,dehghani2018universal,saunshi2025latentthoughts,bay2025mixture,geiping2025scalinglatent,zhu2025scalinglatent,frey2026adaptive,huang2026equilibrium}.

A growing body of work suggests that recurrent computation is especially well suited to complex problems. Looped Transformers can outperform vanilla Transformers on in-context learning and data-fitting tasks, and they can implement multi-step gradient descent in context with far fewer parameters \citep{yang2023looped,fan2024looped,giannou2023looped,gatmiry2024can,gatmiry2024role,chen2025bypassing}. Studies of looped Transformers have also shown that recurrent computation allows shallow parameterizations to approach the performance of deeper untied models on reasoning tasks \citep{saunshi2025latentthoughts,jeddi2026loopformer,geiping2025scalinglatent,zhu2025scalinglatent}. These advantages are particularly visible on abstract reasoning tasks such as ARC-AGI, Sudoku, and Maze, where results from recent recurrent models suggest that recurrence provides a powerful inductive bias for compositional generalization, rule discovery, and inductive reasoning \citep{wang2025hierarchical,jolicoeurmartineau2025trm,gao2025universal}.

Recent work has begun scaling looped Transformers to billion-parameter language models. Ouro models scale to 1.4B and 2.6B parameters with four recurrent steps, reporting strong parameter efficiency relative to larger dense baselines \citep{zhu2025scalinglatent}. Huginn scales to 3.5B parameters and uses 32 recurrent steps to trade parameters for latent computation \citep{geiping2025scalinglatent}. Other systems further study adaptive recurrence, efficient reasoning, memory--compute trade-offs, and scaling laws for stable looped language models \citep{bay2025mixture,jeddi2026loopformer,frey2026adaptive,prairie2026parcae}. However, these gains expose a fundamental compute-accounting issue: \textbf{looping a model $N$ times during pre-training also multiplies pre-training compute by $N$. Thus, looped Transformers should be compared not only against vanilla Transformers with the same parameter count but also against non-looped models trained under the same pre-training compute budget.} This view is consistent with language-model scaling and compute-optimal training analyses, which evaluate quality as a function of parameters, data, and total training FLOPs rather than parameter count alone \citep{kaplan2020scaling,hoffmann2022training,prairie2026parcae}. For example, Ouro-2.6B with 4 loops should be compared against a baseline model with a parameter count close to 2.6B × 4 = 10.4B; likewise, after accounting for pre-training compute, Huginn-3.5B with 32 loops should be compared against a much larger 112B-parameter baseline model. This motivates our central question:

\begin{tcolorbox}[colback=ai2offwhite,colframe=ai2pink,boxrule=0.5pt,arc=3pt,left=6pt,right=6pt,top=5pt,bottom=5pt]
\centering
\emph{Can looped Transformers match or exceed vanilla Transformers \\
under the same pre-training compute budget?}
\end{tcolorbox}

We answer this question with the \textbf{Loopie Series}: two looped MoE LLMs, Loopie-20B-A2B and Loopie-6B-A0.6B, each trained with two loop steps. Our key idea is the \textbf{Loopie Recipe}, a compute-matched scaling recipe that addresses the main concern for recurrent-depth scaling: under a fixed pre-training compute budget, vanilla parameter scaling can otherwise dominate looping \citep{hoffmann2022training,prairie2026parcae}. With a novel post-training method, Loopie develops strong reasoning abilities and achieves frontier-level reasoning performance.

Our contributions are threefold:
\begin{itemize}
    \item \textbf{Compute-matched scaling.} We introduce the Loopie Recipe to address the fixed-compute challenge for looped Transformers and validate it through extensive ablations.
    \item \textbf{Scalable looped MoE models.} We demonstrate that looped computation scales to large MoE language models by training Loopie-20B-A2B and Loopie-6B-A0.6B.
    \item \textbf{Large-scale post-training.} We scale Loopie through large-scale post-training, including a novel supervised pre-training stage, it yields strong reasoning abilities.
\end{itemize}
